\documentclass[10pt,a4paper]{amsart}
\usepackage[margin=19mm]{geometry}
\usepackage{amsmath,amssymb,mathtools}
\usepackage[scr=rsfs,cal=euler]{mathalfa}
\usepackage{xfrac}
\usepackage[unicode,hidelinks,bookmarks=false]{hyperref}
\newcommand{\ie}{i.e.\ }
\newcommand{\cf}{cf.\ }
\newcommand{\resp}{respectively\ }
\newcommand{\Subsection}{\subsection}
\providecommand{\nobreakdash}{\nobreak\hbox{-}}
\setlength{\emergencystretch}{2em}
\multlinegap0pt
\usepackage[all]{xy}
\usepackage{amssymb}
\usepackage{xcolor}
\newtheorem{obs}{Remark}
\newtheorem{teo}{Theorem}
\newtheorem{prop}{Proposition}
\newtheorem{pf}{Pf}
\newcommand{\B}{{\mathcal B}}
\newcommand{\C}{{\mathcal C}}
\newcommand{\Ext}{{\textnormal{Ext}}}
\newcommand{\Hom}{{\textnormal{Hom}}}
\newcommand{\im}{{\textnormal{Im }}}
\title{Projective resolutions of simple modules and Hochschild cohomology for incidence algebras}
\author{Viktor Bekkert, John William MacQuarrie, J\'ulio Marques}
\date{Source: arXiv:2411.07910v2 (CC BY 4.0)}
\begin{document}
\maketitle

\noindent\emph{Source and licence.} Projective resolutions of simple modules and Hochschild cohomology for incidence algebras. Version arXiv:2411.07910v2 (CC BY 4.0), distributed by arXiv under the Creative Commons Attribution 4.0 International licence, http://creativecommons.org/licenses/by/4.0/, as recorded in the arXiv licence metadata for this exact version and shown on the abstract page https://arxiv.org/abs/2411.07910v2. Full licence notice: \texttt{licenses/arxiv-2411.07910-CC-BY-4.0.txt}.\par\smallskip

\noindent\textit{Editorial scope.} Counted unit: the article's prop extgroups (source0.tex lines 354--357). The article's complete original proof of it is delivered with it (source0.tex lines 358--370, one \texttt{proof} environment of 13 lines). No mathematics is written, repaired or re-derived: the statement and the proof are the article's own lines, in the article's own notation, and no retained line is substituted or re-flowed.  Retained uncounted as the source-local context the proof needs: lines 225--227; lines 311--315.  Nothing was omitted from any retained span, so the package carries no omission record.  No retained line contains a citation, so the wrapper carries no bibliography and no bibliography entry.  No retained line contains a figure environment, an image-inclusion command or drawing code, so no image is delivered and the package declares no asset dependency.  Editorial formatting only, in addition: an AMS \texttt{amsart} preamble that restates the article's own declarations and macros used by the retained lines, namely the environments \texttt{obs}, \texttt{teo}, \texttt{prop}, \texttt{pf} on separate counters, together with the standard packages those lines need.  The retained environments print in this document as Obs 1, Teo 1, Proposition 1, Pf 1 in order of appearance, whereas the article prints the same items with the numbering of its own full document.  The complete native source of the article as distributed in the arXiv e-print for this exact version is bundled unchanged as \texttt{sources/168410/source0.tex}.
\begin{obs}\label{obs z' are less than z}
We make explicit a basic but fundamental observation about the $i$-cycles, because it will be used a great deal in what follows.  Namely, if ever $(w,z)\in k\mathcal{C}^{i+1}$ for some $i\geqslant 1$  and $(w',z')\in k\mathcal{C}^i$ is in $supp(w)$ then, $w$ being an element of $(\ker\partial_{i})_z$, necessarily $z'<z$.
\end{obs}

\begin{teo}\label{projectiveresolution}
Let $X$ be a poset and $x\in X$. With definitions as above, the complex
	$$\xymatrix@R=0.8em{ \cdots \ar[r] & k\C^2\otimes_{\Lambda_0}\Lambda\ar[r]^-{d_2} & k\C^{1}\otimes_{\Lambda_0}\Lambda\ar[r]^-{d_{1}} & k\C^{0}\otimes_{\Lambda_0}\Lambda\ar[r] & S_x\ar[r] & 0}$$
where, writing $S_x = \langle s\rangle$, the final non-zero map sends $x\otimes c_{x,a}$ to $\delta_{x,a}s$, is a minimal projective resolution of the simple $\Lambda$-module $S_x$.
\end{teo}

\begin{prop}\label{extgroups}
Let $\Lambda=kX/I$ be an incidence algebra. Let $x,b$ be elements of $X$ and $S_x,S_b$ the corresponding simple $\Lambda$-modules.  For $i\geqslant 0$, we have
	$$\dim_k\Ext_{\Lambda}^i(S_x,S_b) = |\B^{i}_b|.$$
\end{prop}

\begin{pf}
By Theorem \ref{projectiveresolution}, the complex below is a deleted projective resolution of $S_x$.
	$$\xymatrix@R=0.8em{ \cdots \ar[r] & k\C^2\otimes_{\Lambda_0}\Lambda\ar[r]^-{d_2} & k\C^{1}\otimes_{\Lambda_0}\Lambda\ar[r]^-{d_{1}} & k\C^{0}\otimes_{\Lambda_0}\Lambda\ar[r] & 0}$$ 	
Applying the contravariant functor $\Hom_{\Lambda}(-,S_b)$ and denoting by $d^*$ the map $(-\circ d)$ as standard, we obtain the complex
	$$\xymatrix@R=0.8em{ 0 \ar[r]& \Hom_{\Lambda}(k\C^{0}\otimes_{\Lambda_0}\Lambda,S_b) \ar[r]^-{d_{1}^*} & \Hom_{\Lambda}(k\C^{1}\otimes_{\Lambda_0}\Lambda,S_b) \ar[r]^-{d_{2}^*} & \cdots } $$
whose homology groups are $\ker d_{i+1}^*/\im d_{i}^* = \Ext_{\Lambda}^i(S_x,S_b)$, $i\geqslant 0$.  If $x\not\leqslant b$ then, for each $i\geqslant 0$, $\Ext_{\Lambda}^i(S_x,S_b) = 0$ and $|\B^{i}_b|=0$. Suppose now that $x\leqslant b$. We have that $k\C^i\otimes_{\Lambda_0}\Lambda\cong \bigoplus_{(w,z)\in \C^i} P_z$. Therefore, 
	$$\dim\Hom_{\Lambda}(k\C^{i}\otimes_{\Lambda_0}\Lambda,S_b) = |\{(w,z)\in \C^i\,:\,z=b\}| = |\B^i_b|.$$ 
The proposition will thus follow from the observation that the maps $d_i^*$ are zero.  To see this, we first claim that for any $l>0$, an element $\rho$ of $\Hom_{\Lambda}(k\C^{l}\otimes_{\Lambda_0}\Lambda,S_b)$ must send elements of the form $y\otimes \lambda$ with $y\in \C^l$ and $yc_b=0$, to zero: if $y = yc_a\in \mathcal{C}^l$ (some $a$) and $\rho(y\otimes \lambda)\neq 0$, then
$$0 \neq \rho (y\otimes \lambda) = \rho (y\otimes c_a\lambda)c_b = \rho(y\otimes 1)c_a\lambda c_b,$$
implying that $a = b$ as claimed, because non-trivial paths act as $0$ on $S_b$.  Fix $\rho\in \Hom_{\Lambda}(k\C^{i-1}\otimes_{\Lambda_0}\Lambda,S_b)$.  Then 
	$$d_i^*(\rho)((w,b)\otimes 1)=\rho d_i((w,b)\otimes 1) = \rho(w\otimes q_b) = 0$$
because $wc_b=0$ (Remark \ref{obs z' are less than z}).
\end{pf}

\end{document}
